\chapter{QCD 与部分子思想：理论根基}

本书的一切建立在二十世纪七十年代初的三项发现之上：强相互作用在短
距离下变弱；部分子分布服从重整化群演化；理论可以在时空格点上正则化
而不破坏规范不变性。每当库内论文需要说明"为什么格点 QCD 值得做"时，
都会引用这三者。

% ----------------------------------------------------------------------
\section{渐近自由}
\papercard{D.\ J.\ Gross and F.\ Wilczek}{非阿贝尔规范理论的紫外行为}
{Phys.\ Rev.\ Lett.\ \textbf{30}, 1343 (1973)；姊妹篇：H.\ D.\ Politzer,
\emph{强相互作用的可靠微扰结果？},
Phys.\ Rev.\ Lett.\ \textbf{30}, 1346 (1973)}
{无 arXiv（印本前时代）| DOI: 10.1103/PhysRevLett.30.1343 / .1346}
\whyselected{"非阿贝尔规范理论具有渐近自由"这一发现使 QCD 成为可计算
的理论、部分子分布成为有意义的对象，是全书的单点之根。库内在陈述做
格点 QCD 的理由处都会引用它（如 BMW 谱与中子--质子质量差论文，以及
库内那篇手征相变博士论文）。}
\physics{对 $n_f$ 味夸克的 Yang--Mills 理论，单圈 $\beta$ 函数为负，
\begin{equation}
  \beta(g) = -\frac{g^3}{16\pi^2}\Big(11 - \frac{2}{3}n_f\Big) + \mathcal{O}(g^5),
\end{equation}
故耦合在对数上随短距离减小：$Q^2\to\infty$ 时 $\alpha_s(Q^2)\to 0$。
深度非弹性散射因此看到近乎自由的夸克（领头阶 Bjorken 标度无关性），
而大距离下耦合增大——禁闭在此未被证明，但作为同一现象的红外侧面而
变得可信。$\beta_0$ 的符号由胶子自相互作用决定：只有非阿贝尔规范理论
才会反屏蔽。}
\impact{库内的每一个全球 PDF 分析（CT18、NNPDF、CJ、JAM）都用劈裂函数
演化其部分子，而这些劈裂函数的存在性正是由该结果保证的；每篇格点
论文在论证"为什么模拟 QCD"时都援引它。它是整棵树的根。}
\cardend

% ----------------------------------------------------------------------
\section{DGLAP 演化}
\papercard{G.\ Altarelli and G.\ Parisi}{部分子语言中的渐近自由}
{Nucl.\ Phys.\ B \textbf{126}, 298 (1977)}
{无 arXiv（印本前时代）}
\whyselected{DGLAP 方程把标度破坏的算符乘积展开观点翻译成本库每篇
PDF 论文所使用的部分子语言。Ji~2013 开山之作、Radyushkin 与 Orginos
等人的伪分布论文、梯度流 PDF 论文等的参考文献中都出现它（约
$\sim$8/99）。}
\physics{夸克分布不是常数而是依赖标度的幅，服从概率形式的演化方程，
\begin{equation}
  \frac{\partial q(x,\mu)}{\partial\ln\mu^2}
    = \frac{\alpha_s(\mu)}{2\pi}\,(P_{qq}\otimes q + P_{qg}\otimes g),
  \qquad
  (P\otimes q)(x)=\int_x^1\frac{dz}{z}\,P(z)\,q\!\left(\frac{x}{z},\mu\right).
\end{equation}
劈裂函数 $P_{ij}$ 由共线因子化计算；在 Mellin 矩空间它们成为作用于矩
$M_n=\int dx\,x^{n-1}q(x)$ 的矩阵——这正是格点矩计算的语言。DGLAP 是
本书主题领域每张图的"纵轴"：格点上在某个标度 $\mu$ 提取的任何分布都
必须先演化才能与唯象比较。}
\impact{准分布与伪分布的匹配核在短距离极限下复现单圈 $P_{ij}$——这一
自洽检验逐字出现在 Ji~2013、Radyushkin~2017、Balitsky--Morris--Radyushkin
2020 及涂抹准 PDF 论文中。现代三圈劈裂函数（Moch--Vermaseren--Vogt
2004）支撑着库内通用的 NNLO 拟合基准。}
\cardend

% ----------------------------------------------------------------------
\section{格点规范理论}
\papercard{Kenneth G.\ Wilson}{夸克的禁闭}
{Phys.\ Rev.\ D \textbf{10}, 2445--2459 (1974)}
{无 arXiv（印本前时代）| DOI: 10.1103/PhysRevD.10.2445}
\whyselected{格点场论的奠基文献、本库每一项计算的 Conceptual ancestor。
在本语料库内被 BMW 质量论文、AI 采样论文、智能体计算论文与手征相变
博士论文引用（直接引用 $\sim$4/99；实质上普遍）。配套的
\texttt{books/} 收藏中有该论文的英文 LaTeX 全文转排及其中文译本。}
\physics{Wilson 用超立方格点取代连续时空而保持规范不变性严格成立：
自由度是沿链节的规范链变量 $U_\mu(n)\in SU(3)$（平行移动），作用量由
方格（plaquette）构成，
\begin{equation}
  S_W = \frac{\beta}{3}\sum_{n,\mu<\nu}\Big(1-\frac{1}{3}\mathrm{Re\,Tr}\,U_{\mu\nu}(n)\Big),
  \qquad \beta = \frac{2N_c}{g^2}.
\end{equation}
规范不变的观测量是沿闭合路径的链变量乘积之迹——Wilson 圈。面积律
$\langle W(C)\rangle\sim e^{-\sigma A}$ 通过非零弦张力 $\sigma$ 标志
禁闭，静夸克势线性上升。强耦合展开证明 $\beta=0$ 处的面积律；渐近
自由保证 $\beta\to\infty$ 处的库仑律：两个区域共存于同一格点，由蒙特
卡洛（第二章）要探索的 crossover 相连。}
\impact{本库 99 篇论文里的每一个系综、每一个关联函数、每一个流时间，
都因 Wilson 写下这篇文章而存在。他的名字还标记了 Wilson 线——沿路径
排序的链变量指数——它将恰好成为决定部分子物理究竟能否从格点提取的
对象（第五、六章）。}
\cardend
